\section{Robustness: what changes under other choices}\label{sec:robust}

Each result in Section~\ref{sec:validation} is stated under one formulation chosen for legibility: one
theme resolution, flat weighting inside the reference windows, symmetric
five-year windows, estimates conditional on grant, and the 2002--2018
registration cohorts. This section reports what each of those choices costs,
and in each case states the more complicated finding beside the simpler one
the body uses.

\subsection{Representation and weighting}

A scoring is fixed by two choices. The \emph{representation} is what the
language is turned into: a distribution over 50 fitted themes, or a
distribution over the class vocabulary's terms. The \emph{weighting} is
how days inside the reference window count: every day equal, or decaying
by half every two years. Every reference is per-filing, anchored on the
filing's own date (Section~\ref{sec:construct}). Four builds exist:

\begin{center}\small
\begin{tabular}{llrr}
\toprule
Representation & Weighting & sd(lead) & sd(atypicality) \\
\midrule
Themes & flat          & $0.102$ & $0.698$ \\
Themes & half-life 2y  & $0.082$ & $0.698$ \\
Terms  & flat          & $0.155$ & $1.101$ \\
Terms  & half-life 2y  & $0.130$ & $1.102$ \\
\bottomrule
\end{tabular}
\end{center}

\noindent All four are compared on the 6.01 million filings every one of
them scored. The first row is the production scoring.

Changing one choice at a time settles the ordering. Correlations of lead
between builds differing in exactly one respect:

\begin{center}\small
\begin{tabular}{lrr}
\toprule
Choice varied & $r$ (lead) & $r$ (atypicality) \\
\midrule
Weighting, flat vs half-life 2y (themes)  & $0.982$ & $1.000$ \\
Weighting, flat vs half-life 2y (terms)   & $0.970$ & $1.000$ \\
Resolution, 50 vs 200 themes\footnotemark              & $0.636$ & $0.828$ \\
Representation, themes vs terms           & $0.460$ & $0.396$ \\
\bottomrule
\end{tabular}
\end{center}

\footnotetext{The resolution row is computed on the five-year-proof sample
rather than the 6.01 million-filing common sample of the other rows.}
\noindent The weighting is close to irrelevant: whether recent language counts
for more than distant language within the window barely moves the score, and it
does not move atypicality at all to three decimals, which is what the
$45$-degree rotation of Section~\ref{sec:construct} implies --- the long axis is
a property of the filing and the weighting acts almost entirely on the small
residual. Resolution and representation matter
most: quadrupling the number of themes leaves the two builds agreeing on 40\%
of the variance in lead, and themes against terms on less than a quarter. A
robustness check that varies only the weighting is therefore close to a
restatement; one that varies the representation or the resolution is a real
test, which is why the term-scored replication in Appendix~\ref{app:scoring}
and the sweep below are the ones the claims lean on.

\subsection{Theme resolution}

Fifty themes over a 62,168-term vocabulary is roughly one theme per
industry; whether the findings depend on that coarseness is testable by
refitting at higher $T$. Raising the theme count does not add coverage --- the
vocabulary is fixed, and a larger $T$ re-partitions the same words --- but it
changes what counts as sitting far from an industry, and could change any sign.
Re-fitting the model at four resolutions on the same stratified sample and the
same analyzer, and re-scoring the software class under per-filing windows:

\begin{center}\small
\begin{tabular}{rrrrr}
\toprule
Themes & Proof contrast, lead & $t$ & Proof contrast, atypicality & $r$ with $T=50$ \\
\midrule
$50$   & $+3.23$pp & $13.8$ & $-14.01$pp & --- \\
$100$  & $+2.71$pp & $11.6$ & $-13.69$pp & $0.690$ \\
$200$  & $+5.20$pp & $22.2$ & $-13.45$pp & $0.640$ \\
$500$  & $+5.69$pp & $24.3$ & $-14.49$pp & $0.641$ \\
\bottomrule
\end{tabular}
\end{center}

\noindent Three readings, in descending order of how much weight they carry.
The sign and the significance are untouched across a tenfold change in
resolution, and so is atypicality, which stays within a point of $-14$pp
throughout. The magnitude is not: it does not move monotonically in $T$, so
no single resolution can be quoted as the software class's effect, and the
figure this paper reports at
$T = 50$ sits at the low end of the range rather than the high one. A fifth
fit at $T = 1000$ exists and is excluded: its first run shipped
undertrained models through an indentation bug, and the refit is not
persisted. And
agreement with $T = 50$ settles at $0.64$ without decaying as $T$ grows, which
says the resolutions share a fixed core and differ in the remainder rather than
drifting apart as the partition refines.

The corpus-wide contrast is steadier than the single class: at $T = 200$,
scored for all 45 classes, it is $+2.23$pp against $+2.29$pp at $T = 50$. Only
those two resolutions exist corpus-wide, so that is a two-point comparison, but
it is the one the paper's estimates run on, and pooling evidently averages away
much of the per-class resolution sensitivity.

\subsection{Fitting noise at one resolution}

The sweep leaves one question open. Scores at different resolutions agree at
$r = 0.64$ to $0.69$ on lead, and neighbouring resolutions agree no better than
distant ones. That has two readings. The number of themes may genuinely change
what is being measured. Or a theme-model fit may be unstable, with two runs settling in different
local optima, in which case the sweep measures
fitting noise and the resolution has little to do with it.

Holding $T$ fixed and moving only the seed separates the two. A second fit at
$T = 200$ on the same sample, the same vocabulary and the same number of
passes, differing from the first in nothing but its random seed, agrees with
it at $r = 0.79$ on lead and $r = 0.87$ on atypicality across 1.11 million
scored software-class filings. The cross-resolution figures were $0.64$ and
$0.83$. Most of what the sweep showed as resolution sensitivity is therefore
present between two fits that differ only in the seed: for atypicality nearly
all of it, for lead well over half. The resolutions are not measuring
different things; each fit is a noisy draw on the same thing.

The estimate moves less than the scores do. On the identical sample, the
five-year-proof lead contrast is $+5.20$pp under the first seed and $+4.88$pp under
the second (SE $0.23$ each), and the atypicality contrast $-13.45$ and
$-14.38$pp. Two consequences follow. A single fit's lead score carries
measurement error of roughly a fifth of its variance, which attenuates every
coefficient on lead toward zero; correcting for it would enlarge the
estimates, not shrink them. And the spread of
magnitudes across the resolution table above, which does not move
monotonically in $T$, should be read as a band about a third of a point wide
from fitting noise alone rather than as a property of any particular $T$.

The substantive claim does not depend on any of this. The five-year-proof
lead contrast on the 2002--2018 registrations each build scores runs
$+2.29$, $+1.95$, $+3.56$ and $+2.92$pp across the four builds in the
order tabulated above, and the atypicality contrast $-4.39$, $-4.43$,
$-3.63$ and $-3.73$pp. Every sign holds; term scoring returns the larger
lead contrast, so the headline estimate is the conservative
representation.

\subsection{Decayed reference windows}\label{sec:robust_decay}

Inside the reference windows every day counts the same. The alternative
weights recent language up, with a half-life of two years inside the same
1,826-day truncation, so that a term the class adopted last month is more
familiar than one it adopted four years ago. Rescoring the corpus that way
changes the scores little and the estimate somewhat. Flat and decayed lead
correlate at $0.980$ on the 3.10 million gate registrations and atypicality
at $1.000$ to three decimals. On the identical registrations, with quintiles
cut within class and registration year, the top-minus-bottom gate contrast
on lead is $+2.29$pp under the flat window and $+1.95$pp under the decayed
one (SE $0.09$ each); on atypicality it is $-4.39$ and $-4.43$pp. The
decayed window therefore returns a lead penalty about a seventh smaller,
with the same sign, the same precision, and the same profile. The flat
window is kept in the body because it is the simpler statement; a reader who
prefers recent language to count for more should take $+1.95$pp as the
estimate.

\subsection{Window width}\label{app:wsweep}

The window is a parameter of the measure, not a property of it, and the five
years used throughout was inherited from the class-year construction rather
than chosen under the current one. Re-scoring the software class at several
widths under per-filing references, and judging each by the five-year-proof
quintile contrast it produces, gives $+2.02$pp at two years, $+2.20$ at
three, $+3.23$ at five and $+3.37$ at seven, each on an identical sample of
452{,}911 registrations and each with a standard error of $0.23$pp. Signal
rises steeply from two to five years and is flat between five and seven: the
$+1.03$pp gain from three to five is well outside the sampling error, the
$+0.14$pp from five to seven is not. A ten-year window returns $+2.89$pp,
but that figure is not comparable to the others --- requiring both windows
to lie inside the corpus costs 22\% of the scored filings and 13\% of the
gate sample, dropping the earliest and latest registration years, which are
also the years whose contrasts differ most. The apparent decline at ten
cannot be separated from that change in composition and is not evidence
about window width. On the range where the sample is fixed, seven years is
nominally the peak and five is within noise of it while retaining more of
the corpus, so five is defensible rather than optimal. The sweep is on one class
and one outcome.

\subsection{Mixing the two windows}\label{app:wmix}

``Leading'' is relative to a reference, and the two references need not be
the same width. A short past and a long future make lead a measure of
foresight --- whether the filing resembles where the class was heading over
the next seven years more than what it wrote in the last three. A long past
and a short future make it a measure of rupture --- whether the filing has
broken with what the class had established over seven years, judged against
only the next three. Scoring every class on the production representation
and per-filing windows at $W \in \{3, 5, 7\}$ years on each side, and
evaluating all five variants on identical samples --- the gate columns on
the 3.10 million registrations 2002--2018, quintiles within
class$\times$registration-year; the registration columns on the scored
filings 1995--2018, quintiles within class$\times$filing-year:

\begin{center}\small
\begin{tabular}{llrrrr}
\toprule
Variant & Past, future & \multicolumn{2}{c}{First-gate failure} & \multicolumn{2}{c}{Registration} \\
 & (years) & Q5$-$Q1 (pp) & classes $>0$ & Q5$-$Q1 (pp) & Q3 $-$ tails \\
\midrule
Foresight     & 3, 7 & $+1.65$ (0.09) & 31 of 45 & $+1.09$ & $+0.50$ \\
Symmetric     & 3, 3 & $+1.69$ (0.09) & ---      & $+0.79$ & $+0.13$ \\
Symmetric     & 5, 5 & $+2.29$ (0.09) & 33 of 45 & $+0.71$ & $+0.28$ \\
Symmetric     & 7, 7 & $+2.58$ (0.09) & ---      & $+0.48$ & $+0.65$ \\
Past-rupture  & 7, 3 & $+3.00$ (0.09) & 37 of 45 & $+0.01$ & $+0.33$ \\
\bottomrule
\end{tabular}
\end{center}

\noindent The symmetric five-year row is the paper's FE-only estimate in
Table~\ref{tab:gate_lpm}, reproduced here to the decimal, which fixes the
scale. The penalty rises monotonically along the mix, from $+1.65$pp under
foresight to $+3.00$pp under rupture, and rupture is the more pervasive
variant, positive in 37 of 44 classes against 31. In software the two are
$+1.82$ and $+5.07$pp. So the use requirement punishes distance from the
class's established language more than it punishes resemblance to its
coming language; but foresight carries a penalty of its own, more than half
the symmetric one, and the separation is a matter of degree rather than
kind.

Registration is a different matter. On the signed axis the top-to-bottom
contrast is at most a point on a 57\% base and the interior-minus-tails
depth at most two thirds of a point, under every mix. The examiner's
response to lead, on this scoring, is close to nil.

\subsection{Without conditioning on grant}\label{sec:robust_uncond}

Figure~\ref{fig:uncond} repeats the lifecycle outcome analysis without the
registration conditioning, on filings 2013--2015 pooled across all classes
(a window whose registrations land by $\sim$2017 and face at least one
Section~8 gate by the 2025 snapshot; the earliest registrations in the
window also face their year-ten renewal, so panel B mixes one and two
gates). Panel A is the unconditional composite, the probability that a
filing both registers and passes the gate, measured over all filings
including the never-registered. It is nearly flat in \dkl{} because two
opposed gradients cancel: registration rises in lead
($61.3$\% to $63.2$\%) while gate survival falls ($44.2$\% to $40.8$\%);
the caption carries the quintile figures.
Conditioning on grant is what separates them. Panel B shows the conditional
gate outcome on the same filings, and it is U-shaped rather than inverse-U
(depth $-2.11$pp), with a clearly penalised leading tail consistent with the
cleaner single-gate cohort of Figure~\ref{fig:s8pooled}.

\begin{figure}[h]\centering
\includegraphics[width=\textwidth]{results/appendix_unconditional.png}
\caption{Filings 2013--2015 pooled across classes, theme-scored on
per-filing reference windows ($n = 1{,}419{,}678$; gated subset
$n = 763{,}902$). \emph{A:} the unconditional composite P(registered and
passed the five-year proof) over all filings. \emph{B:} P(passed $\mid$ registered
and gated) on the same filings. The unconditional view is nearly flat in
lead (22.4\% to 22.1\% across quintiles); conditioning on grant is what
exposes the gradient.}
\label{fig:uncond}
\end{figure}

\subsection{Toggled filters, windows and parameters}\label{sec:robust_matrix}

Each analysis in this paper and in the companion paper was rerun under its
natural parameter toggles, on the identical data. Nothing
reverses a headline claim; two auxiliaries move.

\begin{center}\small
\begin{tabular}{lll}
\toprule
Variant & Baseline $\to$ variant & Reading \\
\midrule
Gate window 3.5--9.0 / 4.5--8.0 & lead $+2.29 \to +2.29$ / $+2.28$pp & window immaterial \\
Registration cohorts 2002--09 & $+2.29 \to +0.98$pp & the era swing \\
Registration cohorts 2010--18 & $+2.29 \to +3.11$pp & penalty concentrated late \\
Atypicality, both cohort halves & $-4.39 \to -4.55$ / $-4.28$pp & stable \\
Internet pattern, narrow core & goods web gap $0.0 \to +2.0$pp & divergence holds \\
Convergence arrival 1\% / 4\% & early$\to$late $+4.4\to+1.6$ & $+5.9\to+1.8$ / $+3.8\to+2.3$ \\
Surge doubling 1.5$\times$ / 3$\times$ & in-surge net $+1.3$pp & $+1.2$ / $+0.9$pp \\
Wave entry window 4y / 8y & order flat & flat under both \\
Recombination floor 2 / 10 & net $-2.8$pp & $-1.2$ / $-3.6$pp \\
Pair-lead presence 0.05 & net $+2.3$pp & $+3.1$pp \\
Combination measures, class 035 & 009 values & recomb $-3.8$; pair-lead $-1.4$ \\
Ladder, counsel / self debuts & reporting lead $-0.19$pp & $-0.26$ / $-0.03$pp \\
Ladder, 2009--13 / 2014--18 & reporting atyp.\ $+0.11$pp & $0.00$ / $+0.19$pp \\
Excluding duplicate-text filings & lead $+2.29$; atyp.\ $-4.39$pp & $+3.05$; $-5.68$pp \\
Excluding Manual-scale texts only & same & $+2.64$; $-5.05$pp \\
\bottomrule
\end{tabular}
\end{center}

\noindent Three of these change what a reader should carry. The lead
penalty at the gate belongs to the 2010--2018 cohorts; the 2002--2009 half
carries a third of it, which is the era swing the companion paper documents, seen
from the other side. The ladder's gradients belong to counsel-represented
debuts --- among self-filers the reporting gradient is a null --- and its
atypicality gradient to the 2014--2018 half. And the pair-lead penalty,
$+2.3$pp in the software class under every parameter tried there, is
$-2.1$pp in advertising and retail on the corpus-wide construction, and
pooled across 43 classes is zero raw and $+0.5$pp net of lead and
atypicality, so it is reported as a bound for one class rather than a
general fact. All artifacts are under
\texttt{paper/v3/\_eval/}.

\subsection{Lead is a small residual}

Past-facing and future-facing KL are highly correlated however scored. On the
complete class-009 corpus under the production scoring they correlate at
$r = +0.988$ ($n = 1{,}109{,}643$). The lead $L$ is therefore a small residual
of two large, stable quantities: both levels have a standard deviation near
$0.70$ nats about a mean of $1.46$, while $L$ has a standard deviation of
$0.109$ --- about a sixth of either level's spread. Only $1.2\%$ of the two
levels' combined variance survives into their difference. That is what makes $L$ sensitive to
anything that inflates the levels; under term scoring the dominant such
thing is document length, and Section~\ref{sec:construct} quantifies the
coupling under both representations. The consequence for the lead is
direct: the standard deviation
of $L$ rises $2.2$--$2.5$-fold across deciles of $A$ under term scoring
and $1.5$--$1.8$-fold under theme scoring, so a unit of lead means roughly
the same thing everywhere only in the theme representation.
Appendix~\ref{app:representation} shows this as a length surface over the
two axes. Two qualifications. Theme scoring returns a finite value for about 65\% of
filings, so the two columns are not drawn from the same set; re-estimating the
term-scored correlations on exactly the theme-scored subset barely moves them,
which puts the contrast down to the representation rather than the sample.
And ``flat'' means no
monotone length gradient rather than constant: class 035 under theme
scoring is mildly U-shaped in $A$. The two scorings agree only
moderately per filing (Pearson $0.46$, Spearman $0.42$ on the 6.01 million
filings both score), which gives their agreement on the mark-level
patterns evidential force and their disagreement on the firm-level
patterns diagnostic force (Section~\ref{sec:robust_hazard}). Across
$W \in \{3, 5, 7\}$, 6.9--10.4\% of filings flip sign. Per-filing noise
attenuates at every aggregation used: firm-year means, token pools,
theme prevalence trajectories.

Alignment with expert innovation judgment is at most suggestive. Firms on
BCG/MIT ``most innovative'' lists sit $+0.09$ standard deviations above the panel mean
under term scoring ($p = 0.22$), $+0.14$ under theme scoring
($p = 0.07$), and $+0.08$ at $T = 200$ ($p = 0.20$). The matched
sample is 41 firms, which is too small to carry weight in either
direction, and the paper draws no inference from it. Patent counts remain
uncorrelated with \dkl{} on the same panel (Pearson $+0.010$).

\subsection{Term scoring as a cross-check}\label{app:scoring}

The snapshot cross-check on the single one-gate-elapsed cohort
(registrations 2016--2018, terminal status 710 versus 701--705/800)
agrees with the event-dated estimates and is scoring-robust: $-5.4$pp
survival across quintiles under theme scoring and $-5.7$pp under term
scoring on the identical sample (Figure~\ref{fig:s8pooled}). Under term scoring, gate survival
\emph{rises} in the raw KL level ($+6.5$pp past-facing, $+7.6$pp
future-facing): atypical text survives the use requirement better
while leading text survives it worse --- the same two-axis
separation as Section~\ref{sec:bluesea}, at the mark level.

\begin{figure}[h]\centering
\includegraphics[width=\textwidth]{results/topic_outcomes_all.png}
\caption{The two mark-level outcomes under theme and term scoring on
identical samples, pooled across classes, by within-class lead quintile
(1 most lagging, 5 most leading). \emph{Left:} the share of registrations
2016--2018 that survived the five-year proof of continued use at years five to six,
read from terminal status codes, where the two representations agree
closely. \emph{Right:} registration completion
(filings 1995--2018), where they do not: the inverse-U in the signed axis
is present under term scoring (depth $+0.048$) and absent under theme
scoring (depth $-0.001$). The gate result is scoring-robust; the
registration curvature on this axis is not, and the registration analysis
reports only the unsigned version.}
\label{fig:s8pooled}
\end{figure}

\subsection{Functional form}

\begin{figure}[h]\centering
\includegraphics[width=\textwidth]{results/curve_shapes.png}
\caption{Binned profiles at the three gates, with the linear reading the
paper's contrasts assume and the shape selected by AIC where that is not
linear. Forty equal-count bins; the dotted line marks the fitted vertex. The
bottom row is the case that matters: a linear fit through the
public-reporting profile slopes downward while the data turn up at both tails.}
\label{fig:shapes}
\end{figure}

Four shapes were fitted to each profile by weighted least squares on 40
equal-count bins (Figure~\ref{fig:shapes}): linear, quadratic, an exponential
$a + b e^{cx}$, and a free-vertex symmetric form $a + b e^{c|x - x_0|}$ meant
to catch a U directly. Selection is by AIC.

\begin{center}\small
\begin{tabular}{llrrrl}
\toprule
Outcome & Axis & $R^2$ line & $R^2$ curve & $\Delta$AIC & Shape selected \\
\midrule
Registration & atypicality & 0.38 & 0.92 & 77.1 & curve, peak at $+1.04$ \\
Registration & lead & 0.08 & 0.17 & 1.9 & (quadratic; weak) \\
Use-proof & atypicality & 0.74 & 0.74 & 0.0 & linear \\
Use-proof & lead & 0.44 & 0.55 & 4.3 & V, vertex $-1.50$ \\
SEC reporting & atypicality & 0.74 & 0.84 & 15.2 & curve, peak at $+2.17$ \\
SEC reporting & lead & 0.13 & 0.59 & 26.3 & V, vertex $+0.30$ \\
\quad listed only & lead & 0.19 & 0.46 & 12.2 & V, vertex $+0.41$ \\
\quad reporting, not listed & lead & 0.07 & 0.67 & 37.2 & V, vertex $+0.24$ \\
\bottomrule
\end{tabular}
\end{center}

\noindent The winning form is a free-vertex $a + b\,e^{c|x-x_0|}$. On the
atypicality rows it is a genuine curve, with fitted exponents between $0.5$
and $1.2$ --- a hump that peaks about one standard deviation above mean
atypicality at registration and about two above it at SEC reporting. On every
lead row the exponent collapses to the order of $10^{-4}$, which turns
$e^{c|x-x_0|}$ into $1 + c|x-x_0|$: what wins there is a \emph{V}, linear in
distance from a vertex. A V beating a parabola says the turn is sharper than a
quadratic allows, not that anything is multiplying.

Vertices are in standard deviations of the axis across filings.

The two rows for SEC reporting are the outcome Section~\ref{sec:bluesea} and
Table~\ref{tab:edgar} actually run on; the listed and reporting-only splits
below them show the same shape on each half, which is why the split does not
change the reading.

AIC here is computed on 40 binned means, not on the millions of underlying
filings, so it ranks shapes rather than testing them. The $77$ at registration
and the $12$ to $36$ at the public-reporting gates are decisive; the $1.9$ at
registration on lead and the $4.3$ at the five-year proof are weak, and those
two rows are better read as ``no worse than linear'' than as curvature.

The non-monotone profiles on signed lead decompose into parts of
different symmetry (Figure~\ref{fig:ldecomp}). Each ventile's deviation
of the listing rate from the base is split, by sequential counterfactual
prediction, into the part atypicality composition alone predicts, the
part exact-duplicate text adds, a linear lead tilt, and a residual. The
symmetric U is carried about half by atypicality composition --- the
spread of $L$ rises with $A$, so both tails of the lead axis hold
atypical filings and the middle holds typical ones --- and about half by
a residual that is itself symmetric in $|L|$: the same conditional lead
magnitude effect Table~\ref{tab:edgar} estimates at $+0.067$pp per
standard deviation. Duplicate-text composition adds almost nothing once
atypicality is held, and the linear tilt is an order of magnitude smaller
than either symmetric part ($-0.07$pp per nat). The fitted vertices above
follow from this arithmetic: a symmetric profile plus a negative tilt has
its minimum shifted right of zero, and a survival hump its maximum
shifted left, which is where the fits put them ($+0.24$ to $+0.41$, and
$-0.4$, standard deviations).

Excluding text-duplicate filings strengthens rather than weakens every
headline gradient. Of the five-year-proof sample, 47\% of registrations
share their exact normalized description with at least one other filing;
dropping them raises the lead contrast from $+2.29$ to $+3.05$pp and the
atypicality contrast from $-4.39$ to $-5.68$pp, and dropping only the
Manual-scale texts (clusters of 100 or more identical filings) gives
$+2.64$ and $-5.05$pp. Identical texts carry identical scores and
cluster at the typical end of the axis, so they dilute both gradients;
neither is created by them.

\begin{figure}[h]\centering
\includegraphics[width=\textwidth]{results/fig_L_decomposition.png}
\caption{The signed-lead listing profile taken apart. \emph{Left:} each
ventile's deviation of P(owner ever in SEC EDGAR) from the 0.65\% base,
on 2.07 million registered debuts 1995--2018, split into the part
predicted by atypicality composition alone (50 quantile bins of $A$),
the part added by an exact-duplicate-text indicator, a linear tilt in
$L$ fitted to the remainder, and the residual; bars stack from zero,
dots are the actual deviations with 95\% intervals. \emph{Right:} the
same profile on all filings and on filings whose normalized description
text is unique in their class; the central dip survives, and deepens,
without duplicated text.}
\label{fig:ldecomp}
\end{figure}

\subsection{Listing against reporting}

``Appears in SEC records'' mixes two populations: companies whose shares trade, and entities that file with
the Commission without ever listing --- private issuers, funds, subsidiaries
registering securities. Splitting the 19{,}889 matched owner strings on whether
their CIK carries a ticker in the Commission's own company-ticker file gives
10{,}556 listed against 9{,}333 that file financial statements without one.

It does not change the finding. Among registered debut owners whose debut
filing is scored, $0.310$\% later appear as listed companies and $0.285$\% as
reporting non-listed ones, and both profiles are U-shaped in lead with vertices
at $+0.41$ and $+0.24$ standard deviations. The listed profile is slightly
flatter and noisier, which is what the smaller and more selected outcome
predicts. Whatever the signed axis is picking up at this margin, it is not an
artifact of counting non-operating registrants as commercial successes, and it
is not specific to the act of listing either.\footnote{Two limits on the split.
The ticker file is a current snapshot, so a company that listed in 2003 and was
acquired in 2011 appears in the non-listed tier; the tiers are therefore
``currently listed'' rather than ``ever listed,'' which will misclassify older
successes downward. And the base rates here are computed on the appendix's own
panel, which conditions on the debut filing carrying a score and so is smaller
than the body's sample.}

\subsection{Filings are not independent draws}

Some of them share an exact position
with another filing, and outcomes cluster within owner --- the
within-owner correlation is 0.38 for gate survival pooled and is mechanically
1.0 for public listing, since listing is a property of the owner and a
firm with hundreds of marks contributes hundreds of identical successes.
The regression tables are unaffected, because the gate regressions cluster
on normalized owner and the firm-level specifications carry one
observation per owner. The raw quintile contrasts of
the companion paper --- the era, theme, internet
and surge tables --- are not: their two-sample standard errors ignore the
owner clustering, and their $t$-statistics overstate precision by roughly a
factor of two, which the text of those sections now carries where the
margins are close. Binning filings without the correction likewise
overstates how much structure the corpus contains: on a
positional grid, correcting cell errors for the clustering design effect ---
the factor by which within-owner correlation inflates the variance of a cell
mean --- cuts
apparent excess dispersion by roughly a quarter on the gate outcome and
by half on listing. Displays of the vocabulary plane are descriptive
here for that reason, and the near-zero pooled correlation between $A$ and
$L$ reported in Section~\ref{sec:construct} does not survive cell by cell
within class, so per-cell contrasts should not be read as within-class
effects.

